\subsection*{Causal ablation: what a feature is necessary for}
\label{sec:causal}

Annotation and co-activation establish that features correspond to biology, but not that the model
uses them. To test necessity we ablated single features at layer~11: the hidden state was encoded
through the SAE, one feature's activation set to zero, the state decoded and substituted, and the
forward pass continued to the output logits. Disruption is the mean absolute change in output
logits at a set of gene positions, and specificity is that quantity on a gene set divided by the
same quantity on all other positions (Fig.~\ref{fig:causal}, \suptab{causal_arms}).

The design of this test determines what it can show. A feature's annotation is assigned from the
genes it most strongly activates on, so evaluating disruption at those same genes asks whether
ablating a feature affects the genes that define it---a question whose answer is close to
tautological. We therefore partitioned the gene positions of each tested feature into three disjoint
sets: genes that are both in the feature's top-20 activating genes and in its annotated term; genes
of the annotated term that are \emph{not} among its top-20 (an evaluation set independent of the
annotation evidence); and all remaining positions. Each feature additionally received its own null:
a random gene set matched to its held-out set in size and in mean-activation decile. Features were
drawn from three arms---the most richly annotated features, a uniformly random sample of annotated
features, and a uniformly random sample of all alive features---so that the result cannot depend on
how features were selected. All 119 tested features come from K562 cells.

Ablation has a large and reproducible effect on the genes a feature activates on. Median
specificity on the top-20 genes against all other positions is 35.5$\times$ in the richly annotated
arm, 64.3$\times$ and 62.1$\times$ in the two random arms, with over 92\% of features exceeding
2$\times$ in every arm. Restricted to the genes that are both in the top-20 and in the annotated
term, the medians are 37.3$\times$, 56.9$\times$ and 66.9$\times$.

On the independent evaluation set the effect disappears. Median specificity on annotated-term genes
outside the feature's top-20 is \textbf{1.04}, \textbf{0.94} and \textbf{0.97} in the three arms.
The matched random gene sets give 0.89, 0.91 and 0.89, so the held-out sets are statistically
indistinguishable from gene sets chosen at random with the same size and activation level: a
per-feature Wilcoxon signed-rank test of held-out against matched random gives $p = 0.76$ and
$p = 0.34$ in the two unbiased arms, and $p = 0.006$ with a median difference of 0.13 in the richly
annotated arm. Fourteen of the 119 features admit no independent evaluation at all, because every
gene of their annotated term already lies within their own top-20.

The conclusion is narrower than the ablation magnitudes alone suggest. Single-feature ablation
produces a specific, reproducible, gene-level effect, and that effect is confined to the genes on
which the feature fires. It does not extend to the remainder of the biological process the feature
is annotated with. Feature-level interventions therefore reveal genuine computational structure
that component-level ablation misses~\cite{kendiukhov2025systematic}, but the structure they reveal
is the feature's own support rather than the pathway it is labelled with.
